\section{\bt{Introduction}{引言}}
\label{sec:introduction}
\bi{A data agent answers an analytical question by exploring the database: it lists tables, inspects values, tests candidate joins, and settles on SQL over many turns. Agents that work on the same database repeat this exploration, and what they rediscover is knowledge rather than queries: which fact rows a metric aggregates, which records count as one business event, and how tables join (Section~\ref{sec:workload-characterization}). Much of it cannot be recovered from the data at all. Enterprise warehouses name columns in terminology that only domain experts understand~\cite{floratou2024nl2sql}, an o1-preview agent solves only 21.3\% of real enterprise text-to-SQL workflows~\cite{lei2025spider2}, and supplying business knowledge closes much of the gap: expert-annotated evidence raises GPT-4's execution accuracy on BIRD from 34.9\% to 54.9\%~\cite{li2023bird}. Keeping what one agent learned so that later agents can reuse it is therefore a natural design, and systems are converging on it: metric layers let engineers declare definitions by hand~\cite{databricks2026metricviews}, business-intelligence systems derive them from script histories~\cite{weng2025datalab}, and agent-first data systems propose a persistent memory that serves the requests of many agents~\cite{liu2026agents,biswal2026agentsm}. The payoff is large: in our end-to-end study, agents that reuse shared definitions answer \DsCondHoldout\% of held-out questions, against \DsNoShareHoldout\% for agents that explore alone, in \ScenHoldTurnsCondModelA\ instead of \ScenHoldTurnsNoShareModelA\ turns.}
{数据智能体通过探索数据库来回答分析问题：多轮列出表、查看取值、检验候选连接，最后确定 SQL。在同一数据库上工作的智能体反复进行这种探索，而它们重新发现的是知识而不是查询：一个指标聚合哪些事实行，哪些记录算作一次业务事件，表之间如何连接（第~\ref{sec:workload-characterization} 节）。其中很多知识根本无法从数据中恢复。企业数仓的列名使用只有领域专家才懂的术语~\cite{floratou2024nl2sql}；在真实企业 Text-to-SQL 工作流上，基于 o1-preview 的智能体只能完成 21.3\%~\cite{lei2025spider2}；补上业务知识能弥补很大一部分差距：专家标注的证据使 GPT-4 在 BIRD 上的执行正确率从 34.9\% 提高到 54.9\%~\cite{li2023bird}。因此，把一个智能体学到的东西保存下来供后来的智能体复用，是一种自然的设计，系统也正在向它靠拢：指标层让工程师手工声明定义~\cite{databricks2026metricviews}，商业智能系统从脚本历史中推导定义~\cite{weng2025datalab}，面向智能体的数据系统提出为众多智能体的请求服务的持久记忆~\cite{liu2026agents,biswal2026agentsm}。收益很大：在我们的端到端实验中，复用共享定义的智能体答对 \DsCondHoldout\% 的留出题，独自探索的智能体为 \DsNoShareHoldout\%，每题轮数从 \ScenHoldTurnsNoShareModelA\ 降至 \ScenHoldTurnsCondModelA{}。}

\bi{Sharing, however, turns one agent's observation into other agents' premise. Consider store return amount. The SQL an agent learns, \code{SUM(sr\_return\_amt)}, is correct because, on the snapshot it explored, a ticket--item key identifies one completed return, the joins do not fan out, and the return date assigns each return to its period; the SQL states none of this, and the schema enforces none of it. When an ETL change later appends an application-state row for each new return, no table or column changes and the shared SQL still executes, but it counts many returns twice, returning 6,588,699.86 for one period where the correct value is 3,294,349.93. The agent that learned the definition is gone; the agents that reuse it never saw its evidence, have no ground truth, and learn nothing from a successful execution. Such data-only updates are routine: type-2 dimensions add a row per changed member, and late-arriving facts and reloads append rows~\cite{kimball2013toolkit}; even a hand-declared join property in a commercial metric layer ``is not validated at runtime''~\cite{databricks2026metricviews}. A shared memory that only stores and retrieves hands every consumer the same wrong answer.}
{然而，共享把一个智能体的观测变成了其他智能体的前提。以门店退货金额为例。智能体学到的 SQL \code{SUM(sr\_return\_amt)} 之所以正确，是因为在它探索的那个快照上，小票--商品键对应一笔已完成的退货，连接不放大行数，退货日期把每笔退货归入其期间；这些性质 SQL 一条也没写明，表结构一条也不强制。之后一次 ETL 改动开始为每笔新退货追加一条申请状态记录：表和列都没有变，共享的 SQL 照常执行，却把许多退货计算了两次，某个期间返回 6,588,699.86，而正确值是 3,294,349.93。学到这个定义的智能体已经结束；复用它的智能体从未见过它的证据，没有基准真值，从一次成功的执行中也看不出任何问题。这类纯数据更新是日常操作：第 2 类缓慢变化维为每个变化的成员追加一行，迟到事实与重新装载追加行~\cite{kimball2013toolkit}；即使是商用指标层中人工声明的连接性质，也“不在运行时验证”~\cite{databricks2026metricviews}。只负责存储和检索的共享记忆，会把同一个错误答案交给每一个使用者。}

\bi{This memory differs from the conversational, factual, and procedural memories that agent systems keep~\cite{chhikara2025mem0,wang2025awm,zhang2026ace} in three ways. Its content describes an external database that keeps changing, so it can become wrong without anyone editing it. Its producers and consumers are different agents, so the consumer cannot vouch for what it reads: in our traces only 16\% of sessions check key uniqueness, and every wrong answer came from SQL that ran without error (Section~\ref{sec:workload-characterization}). And its validity is partly decidable by queries over the same database, so the layer that holds it can check once, for all its consumers, whether it still holds and enforce the outcome where it is used. Existing safeguards cover parts of this. Invalidation on schema change misses data-only breaks; data-validation pipelines~\cite{schelter2018deequ,breck2019validation,dbt2026tests} test hand-declared constraints per table without knowing which shared knowledge depends on them; agent memories retrieve by similarity and leave the consistency of shared memory open~\cite{biswal2026agentsm,liu2026agents}. Telling agents to verify validity themselves recovers little: under the breaking changes of our study, it lifts the accuracy of example retrieval only from \DsTrajModeled\% to \DsTrajVerifyModeled\%.}
{这种记忆与智能体系统保存的对话、事实和流程记忆~\cite{chhikara2025mem0,wang2025awm,zhang2026ace}有三点不同。它的内容描述一个不断变化的外部数据库，因此不需要任何人修改它就可能变错。它的生产者与使用者是不同的智能体，因此使用者无法为自己读到的内容担保：在我们的轨迹中，只有 16\% 的会话检查键唯一性，所有错误答案都来自执行成功、没有报错的 SQL（第~\ref{sec:workload-characterization} 节）。它的有效性又可以部分地由同一数据库上的查询判定，因此保存它的那一层能够替所有使用者统一检查它是否仍然成立，并在使用处执行检查的结论。现有保障手段只覆盖其中一部分。只在模式变更时失效会漏掉纯数据破坏；数据验证管道~\cite{schelter2018deequ,breck2019validation,dbt2026tests}按表测试人工声明的约束，却不知道哪些共享知识依赖它们；智能体记忆按相似度检索，把共享记忆的一致性留作开放问题~\cite{biswal2026agentsm,liu2026agents}。让智能体自己核验有效性收效甚微：在我们实验的破坏性变化下，它只把示例检索的正确率从 \DsTrajModeled\% 提高到 \DsTrajVerifyModeled\%。}

\bi{We therefore build the shared memory as a layer that takes over what each consumer would otherwise have to check alone, and the layer must answer three questions together. \emph{What may be published?} A successful task mixes reusable business meaning with question parameters, incidental observations, and dead ends; the layer must publish only the reusable part, with the evidence and conditions that justify it, and it must check what other agents will read rather than what the producer ran. \emph{When may a use rely on it?} Retrieval decides whether knowledge is relevant, not whether it holds on the data a query reads. A passed check describes one snapshot; a write can commit between the check and the query that relies on it, and cheap version signals such as DML statistics are neither transactional nor immediate: under concurrent writers, check-then-execute answers on condition-violating data in up to \SnPreUnannMax\% of uses. \emph{Who keeps it valid, and who improves it?} If every consumer revalidated what it reads, sharing would save little, so maintenance must be paid once for all consumers; yet both ways of changing shared knowledge are risky. A repair can restore one row per business key while selecting the wrong population, and the only regression reference a deployment has, the producer's SQL, is valid only on the snapshot it was written for; a rewrite that is equivalent on today's data may be equivalent only because of a property that tomorrow's data breaks.}
{因此，我们把共享记忆做成一层，由它承担原本要由每个使用者各自检查的部分；这一层必须同时回答三个问题。\emph{什么可以发布？}一次成功的任务中混杂着可复用的业务含义、题目参数、偶然的观测和走不通的尝试；这一层只能发布可复用的部分，连同支撑它的证据和条件，并且要检查其他智能体将要读到的内容，而不是生产者当时执行的内容。\emph{一次使用何时可以依赖它？}检索判断的是知识是否相关，而不是它在查询所读的数据上是否成立。一次通过的检查只描述一个快照；写入可能在检查与依赖它的查询之间提交，而 DML 统计这类廉价的版本信号既不随事务提交，也不即时可见：存在并发写入时，先检查后执行会让最多 \SnPreUnannMax\% 的使用在违反条件的数据上作答。\emph{由谁保持它有效，由谁改进它？}如果每个使用者都要重新验证自己读到的内容，共享就省不下多少工作，所以维护必须为所有使用者只付一次；而改变共享知识的两种方式都有风险。修复可以恢复每个业务键一行，却选中错误的总体，而部署中唯一的回归参照，即生产者写的 SQL，只在它所针对的快照上成立；在今天的数据上等价的改写，可能只是因为某个明天的数据会破坏的性质才等价。}

\bi{\system\ (\textbf{M}etric-\textbf{A}ware \textbf{V}alidation and \textbf{R}euse for \textbf{A}gents) answers these questions with one abstraction, a versioned \emph{revision} of a definition that carries the executable conditions on which its SQL depends. \emph{Publication}: admission turns a judged-successful trajectory into a structured definition that separates definition filters from question parameters, checks it statically and on the current snapshot, and replays the canonical SQL compiled from its fields. The conditions follow from that structure, namely filtered key uniqueness for the grain, cardinality and completeness for each join, and the role of the date join, and they suffice: together with the definition's business premises, they make the canonical SQL count every business event exactly once in its period, up to the omissions that the completeness conditions bound (Proposition~\ref{prop:sufficiency}). \emph{Reliance}: an agent relies on a revision by declaring it in \code{run\_sql}; \system\ runs the query in a repeatable-read snapshot on which the revision's conditions hold, reusing check results whose transactional versions match the snapshot and checking the rest inside it (Lemma~\ref{lem:reuse}). \emph{Maintenance and improvement}: a check result is keyed by the condition's canonical form and the versions of the tables it reads, so it serves every definition and agent that needs the same condition until one of those tables changes; a failed condition invalidates the revision, and repair publishes a replacement only if it is the unique structural fix and reproduces the producer's answer in an as-of query at learning time (Algorithm~\ref{alg:repair}). The same channel lets the memory improve: \system\ rewrites published definitions, by rule or by LLM, and publishes a rewrite as the next revision only if it returns the current revision's results on every sampled period and as of learning time and is measurably cheaper; a premise that the rewrite adds, such as date keys that are contiguous across months, becomes one more maintained condition, which for the rewrite rules implies equivalence on every period (Proposition~\ref{prop:rewrite}); if a change breaks it, the predecessor is reinstated. Consumers learn of every new, repaired, optimized, or invalidated revision in their next tool response. The managed knowledge is metric definitions with their grain and join paths, and the guarantee covers declared uses; because maintenance operates on structured definitions, it applies equally to definitions declared in a semantic layer.}
{\system（\textbf{M}etric-\textbf{A}ware \textbf{V}alidation and \textbf{R}euse for \textbf{A}gents，面向智能体的指标感知验证与复用）用一个抽象回答这些问题：定义的版本化\emph{修订}，它带着其 SQL 所依赖的可执行条件。\emph{发布}：准入把一条判题成功的轨迹转成结构化定义，把定义过滤与题目参数分开，做静态检查和当前快照上的检查，并重放由字段编译出的规范 SQL。条件由这一结构推出，即粒度对应过滤后键唯一性，每个连接对应连接基数与连接完整性，日期连接另有角色条件；它们是充分的：与定义的业务前提一起成立时，规范 SQL 把每个业务事件在其期间内恰好计数一次，遗漏仅限于完整性条件所约束的部分（命题~\ref{prop:sufficiency}）。\emph{依赖}：智能体在 \code{run\_sql} 中声明一个修订，即表示依赖它；\system\ 在一个该修订条件成立的可重复读快照中执行查询，事务性版本与快照一致的检查结果直接复用，其余条件在快照内检查（引理~\ref{lem:reuse}）。\emph{维护与改进}：检查结果按条件的规范形式及其读取表的版本索引，因此在这些表变化之前，它服务于所有需要同一条件的定义和智能体；条件失败使修订失效，只有当替代修订是结构上唯一的修复、且在学习时刻的 as-of 查询中复现生产者的答案时，修复才发布它（算法~\ref{alg:repair}）。同一通道也让记忆得以改进：\system\ 用规则或大模型改写已发布的定义，只有当改写在每个抽样期间上、以及按学习时刻都返回与当前修订相同的结果，并且可测量地更省时，才把它发布为下一个修订；改写新增的前提（例如日期键跨月连续）成为又一条受维护的条件；对改写规则，它蕴含改写在每个期间上的等价（命题~\ref{prop:rewrite}），数据变化破坏它时恢复前一修订。使用者在下一次工具调用的返回中得知每一个新发布、已修复、已优化或已失效的修订。受管理的知识是指标定义及其粒度与连接路径，保证覆盖声明了修订的使用；由于维护作用于结构化定义，它同样适用于在语义层中声明的定义。}

\bi{We evaluate \system\ at the system level, where what agents learned is fixed and only maintenance varies, and end to end with DeepSeek V4.1 Flash agents. Replaying \RpLibs\ libraries learned by three LLMs and a library derived from the 99 TPC-DS query templates under \RpChanges\ data changes, maintained definitions serve no wrong answer under the covered changes, whereas invalidation on schema change serves \RpSchemaWrongModeled\ and \TrSchemaWrongModeled. Same-snapshot use answers on condition-violating data in none of \SnSnapAnswered\ uses. Regression as of learning time repairs as much as a gold reference does, and without the uniqueness rule an indistinguishable backup copy yields a wrong published repair. Maintenance work follows the distinct conditions an update touches, not the definitions and agents that depend on them: with 19 shared definitions it costs \CbStagCondSix\% of a full recheck per definition, a saving that a version-keyed cache of the check queries also obtains. End to end, sharing makes held-out tasks \AmHoldSpeedup$\times$ faster with \AmHoldTokSave\% fewer tokens, and maintained memory answers \DsCondModeled\% under the covered breaking changes, against \DsTrajModeled\% for a matched example-retrieval baseline that is as accurate on unchanged data; admission refused candidates whose canonical SQL would have returned up to \PubExRatio\ times the learned answer; and a rule rewrite made \OptPublished\ of \OptDefs\ learned definitions \OptSaveMin--\OptSaveMax\% faster, falling back to the predecessors without a wrong answer when a change broke its premise.}
{我们在系统层（固定智能体学到的内容，只改变维护方式）和端到端层（基于 DeepSeek V4.1 Flash 的智能体）评估 \system。把 3 个大模型学到的 \RpLibs\ 个定义库，以及从 TPC-DS 的 99 个查询模板导出的定义库，在 \RpChanges\ 种数据变化下回放：在条件覆盖的变化下，维护后的定义不给出错误答案，只在模式变更时失效的做法分别给出 \RpSchemaWrongModeled\ 和 \TrSchemaWrongModeled\ 个。同快照使用在 \SnSnapAnswered\ 次使用中没有一次在违反条件的数据上作答。按学习时刻作回归测试，修复的题数与标准答案参照相同；没有唯一性规则时，无法区分的备份副本会导致发布错误的修复。维护工作沿更新触及的不同条件展开，而不是沿依赖它们的定义和智能体展开：19 个定义共享条件时，维护耗时为整定义重查的 \CbStagCondSix\%，按版本索引检查查询的缓存也能得到这一节省。端到端实验中，共享使留出题任务快 \AmHoldSpeedup\ 倍、token 少 \AmHoldTokSave\%；在条件覆盖的破坏性变化下，维护后的记忆答对 \DsCondModeled\%，数据不变时同样准确的匹配的示例检索基线为 \DsTrajModeled\%；准入拒绝了规范 SQL 会返回最多为所学答案 \PubExRatio\ 倍的候选；一条改写规则让 \OptDefs\ 条学到的定义中的 \OptPublished\ 条快了 \OptSaveMin\%--\OptSaveMax\%，数据变化破坏其前提时回退到前一修订，没有给出错误答案。}

\bi{This paper makes the following contributions:}
{本文的贡献如下：}
\begin{itemize}
\item \bi{We identify what a shared memory layer for data agents should take over so that no consumer carries it alone, namely publication, reliance, and maintenance and improvement, and realize it with versioned revisions of metric definitions that carry four kinds of typed, executable validity conditions, which we prove sufficient for the canonical SQL to count every business event exactly once in its period, up to omissions bounded by the completeness conditions (Sections~\ref{sec:overview} and~\ref{sec:model}).}
{我们指出数据智能体的共享记忆层应当接过哪些不应由每个使用者各自承担的职责，即发布、依赖以及维护与改进，并用指标定义的版本化修订实现它；修订带有四类类型化、可执行的有效性条件，我们证明这些条件足以让规范 SQL 把每个业务事件在其期间内恰好计数一次，遗漏仅限于完整性条件所约束的部分（第~\ref{sec:overview} 和~\ref{sec:model} 节）。}
\item \bi{We bind every declared use to validity evidence that holds on the query's own snapshot: check results keyed by canonical form and dependency versions, version checks and notices, and same-snapshot validation, which applies the transactional-consistency principle of TxCache~\cite{ports2010txcache} to check results under a reuse rule we prove sound (Lemma~\ref{lem:reuse}; Section~\ref{sec:execution}).}
{我们把每一次声明了修订的使用绑定到在查询自身快照上成立的有效性证据：按规范形式与依赖版本索引的检查结果、版本校验与通知，以及同快照验证；后者把 TxCache~\cite{ports2010txcache} 的事务一致性原则用于检查结果，其复用规则经我们证明是可靠的（引理~\ref{lem:reuse}；第~\ref{sec:execution} 节）。}
\item \bi{We make maintenance and improvement a shared service: condition-level revalidation shared across definitions and agents, and a bounded repair algorithm that publishes only the unique fix that reproduces the producer's answer as of learning time (Algorithm~\ref{alg:repair}); the same revision channel carries improvements, adopting a cheaper rewrite whose equivalence the maintained conditions imply (Proposition~\ref{prop:rewrite}) and reinstating its predecessor when a change breaks the added premise (Section~\ref{sec:maintenance}).}
{我们把维护与改进做成共享的服务：在定义与智能体之间共享的条件级重验证，以及有界修复算法，只发布按学习时刻复现生产者答案的唯一修复（算法~\ref{alg:repair}）；同一修订通道也承载改进：只采纳受维护条件蕴含其等价的更省改写（命题~\ref{prop:rewrite}），数据变化破坏新增前提时恢复前一修订（第~\ref{sec:maintenance} 节）。}
\item \bi{We implement \system\ and evaluate it at the system level and end to end, separating the value of sharing from the value of maintenance and measuring what building the memory costs and how much faster consumers become. A characterization of data-agent workloads against 441 million queries from 400 Amazon Redshift clusters motivates the design: agents rediscover knowledge that no query cache captures, rarely check validity themselves, and fail silently (Sections~\ref{sec:background} and~\ref{sec:evaluation}).}
{我们实现了 \system，并在系统层和端到端层评估它，把共享的价值与维护的价值分开测量，并测量建立记忆的代价以及使用者快了多少。把数据智能体负载与 400 个 Amazon Redshift 集群的 4.41 亿条查询对照的负载刻画支撑了这一设计：智能体反复重新获取查询缓存捕获不到的知识，很少自行检查有效性，并且静默出错（第~\ref{sec:background} 和~\ref{sec:evaluation} 节）。}
\end{itemize}
